\scanpage{296}
\mspage{2\unsure{0|9}}
% archive sheet number: 295

Mais ici me vient l'idée du lemme de K.T. Brown, cité par \unsure{Baues}, que je vais redonner ici :

\uline{Lemme} 3.2.1. (K.T. Brown) Soit $(M,W)$ un localiseur ($W$ modérément saturé), muni d'un ens. $C\subset\mathrm{Fl}\,M$ satisfaisant les conditions :

1) $C$ stable par composition et par cochangement de base (NB on n'a \deleted{pas} à supposer l'existence des sommes amalgamées
\diagram{\begin{tikzcd}[ampersand replacement=\&,row sep=small,column sep=small]
X_0\arrow[rr,"i_1"]\arrow[dr,"i_2"']\& \&X_1\arrow[dr] \&\\
\& X_2\arrow[rr,no head] \& \&X_3
\end{tikzcd}}
que si $i_1\in C$), contient les iso.

2) $M$ contient un objet initial, $\varnothing_M$, et pour tout objet \inserted{$X$} de $M$, $\varnothing\longrightarrow X$ soit dans $C$. (NB 1) et 2) impliquent que $M$ est stable par sommes finies)

% A long arrow carries this item above the preceding condition.
\inserted{\deleted{\unsure{4}}3) Toute flèche $f$ de $M$ se factorise en $pi$, avec $i\in C$, $p\in W$.}

\deleted{3}4) $C\cap W$ est stable par cochangement de base, i.e. si
\diagram{\begin{tikzcd}[ampersand replacement=\&,row sep=small,column sep=small]
X_0\arrow[rr,no head,"i_1"]\arrow[dr,"{i_2\deleted{\unsure{i_2}}}"'] \& \&X_1\arrow[dr,"{\deleted{\unsure{j_1}}j_2}"] \&\\
\& X_2\arrow[rr,no head,"j_1"'] \& \&X_3
\end{tikzcd}}
est cocartésien et \deleted{\unsure{ons}} $i_1\in C\cap W$, alors $j_1\in C\cap W$. (NB $j_1\in C$ est déjà contenu dans 1), il reste \inserted{à exiger} $j_1\in W$)

Sous ces conditions, on a $\deleted{\unsure{W}}C\subset\mathrm{Cof}_{\unsure{W|w}}$.

On peut aussi énoncer le lemme ainsi

\uline{Cor.} 3.2.2. Supposons que $C$ satisfait 1, 2, 3. Alors \deleted{la condition 3) \uline{équivaut} à $C\subset\mathrm{Cof}_W$,}
